\chapter{Conclusion}

\section{Summary of Findings}

% To be completed after final results review.

\section{Contributions and Impact}

% To be completed.

\section{Limitations of the Study}

% To be completed.

\section{Future Work}

Heading level detection was the weakest metric across all four strategies,
with no strategy exceeding 0.565.
The current approach infers heading levels from font-size ratios for the
text path and from visual appearance for the vision path.
Neither method reliably assigns the correct depth when the document uses
inconsistent font sizes or non-standard heading styles.
A dedicated heading classifier trained on labelled thesis documents would
likely improve this score more than any other single change.

\begin{itemize}

    \item \textbf{Table complexity classifier.}
    The page classifier currently routes every page with a detected table
    to the vision LLM.
    A refinement that distinguishes simple row-column tables from complex
    ones --- merged cells, irregular spans, embedded figures --- would
    keep simple tables in the faster text path without sacrificing quality
    on the pages where the VLM actually adds value.

    \item \textbf{Dedicated formula recognition.}
    Formula pages currently use the general-purpose Qwen2.5-VL-7B model.
    A specialised tool such as pix2tex, which is trained exclusively on
    mathematical notation, would produce cleaner LaTeX output and run
    faster on formula-heavy pages.
    The adaptive architecture supports this substitution without changes
    to the rest of the pipeline.

    \item \textbf{Larger and domain-specific models.}
    All VLM results in this experiment use Qwen2.5-VL-7B.
    A larger model or one fine-tuned on academic documents could reduce
    the paraphrasing tendency that causes image-only and hybrid to score
    poorly on paragraph and code block structure.
    Retesting with a different model would show whether the strategy
    rankings in Chapter~6 are model-dependent or consistent.

    \item \textbf{Broader document corpus.}
    The evaluation covers three German bachelor theses.
    Testing on research papers, lecture slides, and scanned documents
    would reveal whether the page classifier and postprocessing pipeline
    generalise beyond the thesis format.
    Scanned documents in particular would require OCR integration, which
    the current system does not provide.

    \item \textbf{Automated output comparison tool.}
    Manual inspection of converted pages, as done when analysing the
    HITL thesis output, is slow and error-prone.
    A script that shows a per-page diff between converted output and
    reference files, highlighting only pages below a similarity
    threshold, would make iterative improvement of the postprocessing
    pipeline significantly faster.

\end{itemize}
